% Draft prepared 2026-10-11; elementary exact observations, no novelty claim.
% Author: analytic_transfer agent. Independent review is requested separately.
\begin{proposition}[The moment-generating-function boundary of the clause hitting time]
\label{prop:hitting-mgf-boundary}
Fix integers $n\ge k\ge2$, and let $C_1,C_2,\ldots$ be independent
uniform proper $k$-clauses on $n$ variables. Write
$F_m=\bigwedge_{i=1}^m C_i$, with $F_0=\mathrm{True}$, and let
\[
 T=\min\{m\ge0:F_m\text{ is UNSAT}\},\qquad q=1-2^{-k}.
\]
Then, for every integer $m\ge0$,
\begin{equation}
 q^m\le \mathbb P(T>m)\le 2^nq^m.
 \label{eq:hitting-tail-sandwich}
\end{equation}
In particular, $T$ is finite almost surely, and its moment generating
function is finite precisely for real
\[
 \theta< -\log(1-2^{-k}).
\]
For $0<\theta< -\log q$, the explicit bounds are
\[
 1+\frac{e^\theta-1}{1-qe^\theta}
 \;\le\;\mathbb E e^{\theta T}
 \;\le\;1+\frac{2^n(e^\theta-1)}{1-qe^\theta}.
\]
\end{proposition}
\begin{proof}
Fix one assignment. Each clause is satisfied by that assignment with
probability $q$, independently, so the probability that it satisfies
$F_m$ is $q^m$. This event implies $T>m$. Conversely, if $T>m$, some
one of the $2^n$ assignments satisfies $F_m$. A union bound establishes
\eqref{eq:hitting-tail-sandwich}. Its upper bound tends to zero as
$m\to\infty$, proving almost-sure finiteness.

For $\theta>0$, the pointwise geometric-sum identity and Tonelli's
theorem for nonnegative summands give, allowing infinite expectations,
\[
 \mathbb E e^{\theta T}
 =1+(e^\theta-1)\sum_{m=0}^{\infty}
                  e^{\theta m}\mathbb P(T>m).
\]
The two bounds in \eqref{eq:hitting-tail-sandwich} show that this series
converges if and only if $qe^\theta<1$, and give the displayed estimates
when it converges. At equality it diverges by the lower bound. For
$\theta\le0$, finiteness follows from $0<e^{\theta T}\le1$.
\end{proof}

\begin{example}[The unmodified Poisson variance bound fails]
\label{ex:coupon-hitting-time}
When $n=k=2$, each of the four possible clauses excludes exactly one
of the four assignments. Thus $T$ is the time until all four clause
types have appeared. Successive waiting times for a new type are
independent geometric random variables, supported on
$\{1,2,\ldots\}$, with success probabilities
$1,3/4,1/2,1/4$. Their means and variances yield
\[
 \mathbb E T=1+\frac43+2+4=\frac{25}{3},\qquad
 \operatorname{Var}(T)=0+\frac49+2+12=\frac{130}{9}
 >\mathbb E T.
\]
Independence follows from independence of the future draws at each
collection time and from the fact that the probability of a new type
depends only on the number already collected.

Consequently, the rotation-volume inequality
$\operatorname{Var}(I)\le\mathbb E I$ does not transfer verbatim to
$T$. More generally, Proposition~\ref{prop:hitting-mgf-boundary}
rules out any finite all-real Poisson cgf bound for $T-\mathbb E T$,
such as $\log\mathbb E e^{\theta(T-\mathbb E T)}
\le a(e^\theta-\theta-1)$ with finite $a$ for every real $\theta$.
These obstructions leave open estimates on a restricted tilt interval
with suitable $k$-dependent constants.
\end{example}
